\noindent
This paper asks how a promise that carries no enforcing force\,---\,a vow, a betrothal document, a memorandum of understanding\,---\,can nonetheless generate real and well-founded trust. Beginning from the symptomatology of the contemporary condition, in which the apparatus of assurance proliferates at every scale while the trust it secures grows scarcer, it diagnoses the failure to a mistake about causation: the taking of force, the mechanical cause, for the cause of an effect, trust, that is inferential and free. It traces, through the philosophy of law, the migration of the anchor of legal validity from force toward trust, which the law records at its own technical edge in the doctrine of promissory estoppel; and it argues that the inferential cause is a different kind of thing from the mechanical, the kind that addresses a free being as free. Its central contribution is a mechanism: trust is the precision a receiver's generative model assigns to the prediction that a promise will be kept, generated through three coupled channels\,---\,a structural prior, a high-fidelity self-committing likelihood, and a wagered prior\,---\,whose multiplicative coupling makes genuine trust robust where any single channel is forgeable. It gives this mechanism a geometric reading in the language of holonomy and a structural reading, for relational trust, in the language of open quantum-system dynamics; states the criterion of justice that distinguishes genuine trust from forged; develops a theory of the translation of sincerity across incommensurable epistemologies; and reconceives, in a historical-materialist register, the optimisation of the juridical as the decentralisation of enforcement. It rises, at its summit, to the claim that the self-legislation, self-adjudication, and self-keeping of a subject who binds himself without recourse is the descent of a relational divinity\,---\,so that the non-binding vow proves, in the dimension of the divine, the purer realisation of what a promise is. It concludes in plural conclusions, refusing, in its own form, the synthesis its thesis forbids.
